\documentclass[11pt]{article}
\usepackage[T1]{fontenc}
\usepackage[utf8]{inputenc}
\usepackage{mathptmx}
\usepackage[margin=1in]{geometry}
\usepackage[hidelinks]{hyperref}
\usepackage{enumitem}
\usepackage{microtype}
\usepackage{tabularx}
\usepackage[dvipsnames]{xcolor}
\definecolor{ReviewerGreen}{HTML}{008000}
\newcommand{\rvcomment}[1]{{\color{ReviewerGreen}#1}}
\newcommand{\concern}[1]{\textbf{\color{ReviewerGreen}#1}}
\setlength{\parindent}{0pt}
\setlength{\parskip}{6pt}
\setlist[itemize,1]{leftmargin=1.6em,labelsep=0.6em,itemsep=3pt,topsep=3pt,parsep=0pt,
                    label=\raisebox{-0.15ex}{\large\textbullet}}
\setlist[itemize,2]{leftmargin=1.4em,labelsep=0.5em,itemsep=2pt,topsep=3pt,parsep=0pt,
                    label=\textendash}
\pagestyle{plain}
\begin{document}
\noindent\begin{tabularx}{\textwidth}{@{}l@{\hspace{1.1em}}X@{}}
\textbf{Original Manuscript ID:} & Access-2026-39991 \\[3pt]
\textbf{Original Article Title:} & ``Version Migration as a Correlated Shock: Model Updates as an Unmanaged Risk Channel in LLM-Based Financial Systems''
\end{tabularx}

\vspace{18pt}

\noindent\begin{tabularx}{\textwidth}{@{}l@{\hspace{1.1em}}X@{}}
\textbf{To:} & IEEE Access Editor \\[3pt]
\textbf{Re:} & Response to reviewers
\end{tabularx}

\vspace{16pt}
Dear Editor,

Thank you for allowing a resubmission of our manuscript, with an opportunity to address the reviewers' comments.

We are uploading (a) our point-by-point response to the comments (below) (response to reviewers, under ``Author's Response Files''), (b) an updated manuscript with yellow highlighting indicating changes (as ``Highlighted PDF''), and (c) a clean updated manuscript without highlights (``Main Manuscript'').

\vspace{14pt}
Best regards, \\[22pt]
Samir Chincholikar \\
Robin Chawla

\newpage

\vspace{12pt}
\concern{Reviewer \#1, Concern \#1}

\concern{Reviewer comment:}

\rvcomment{1. Redefine or defend the excess-flip estimand}

\rvcomment{The within-version rates compare different replicate seeds, whereas Fcross compares the same nominal seed across versions. Because seed behavior is provider and version specific, these quantities need not share a common coupling.}

\textbf{Author response:}

\begin{itemize}

\item We agree that Fwithin and Fcross use different constructions and need not share a common coupling.

\item We therefore no longer interpret Fexcess as a variance-reduced within-cell contrast or as the causal share of decisions changed by migration.

\item As an alternative sensitivity check, we estimated same-seed repeat-call disagreement on a stratified subgrid.

\item The repeat-call floors are 0.028 for OpenAI and 0.014 for Gemini.

\item Applying these alternative noise-floor estimates to the full-grid cross-version disagreement rates gives excess disagreement of 0.189 [0.147, 0.206] for OpenAI and 0.375 [0.339, 0.385] for Gemini.

\item The substantive conclusion is unchanged.

\end{itemize}

\textbf{Author action:}

\begin{itemize}

\item Section III-C, ``What the Two Baselines Measure, and a Repeat-Call Sensitivity,'' was added.

\item The section now:

\begin{itemize}

\item explicitly distinguishes the construction of Fwithin and Fcross;

\item narrows the interpretation of Fexcess;

\item reports a same-seed repeat-call sensitivity analysis on a stratified subgrid;

\item clearly states that the resulting repeat-call floor is applied as an alternative noise-floor estimate to the full-grid cross-version disagreement rate; and

\item retains the preregistered endpoint as primary.

\end{itemize}

\end{itemize}

\vspace{12pt}
\concern{Reviewer \#1, Concern \#2}

\concern{Reviewer comment:}

\rvcomment{2. Separate distributional distance from per-call migration flips:}

\rvcomment{If estimated with independent replicate draws, the proposed difference has a defensible categorical distribution-distance interpretation. That is different from the probability that a specific production decision will flip after cutover. Explain which operational question each statistic answers and avoid treating them as interchangeable.}

\textbf{Author response:}

\begin{itemize}

\item We agree that these quantities answer different operational questions.

\item TVD measures movement in the aggregate label distribution.

\item Fcross measures matched decision disagreement across versions.

\item We also clarify that TVD equals the marginal component M of the prevalence-churn decomposition.

\end{itemize}

\textbf{Author action:}

\begin{itemize}

\item Section III-D, ``Distribution Distance Is Not Per-Call Flip Probability'' was added.

\item The new section:

\begin{itemize}

\item formally distinguishes TVD from Fcross;

\item reports both quantities side by side for both migration pairs;

\item explains their different operational interpretations; and

\item links TVD explicitly to the marginal component M.

\end{itemize}

\end{itemize}

\vspace{12pt}
\concern{Reviewer \#1, Concern \#3}

\concern{Reviewer comment:}

\rvcomment{5) Are there references that are not appropriate for the topic being discussed?: Yes}

\rvcomment{5a) If yes, then please indicate which references should be removed.: Reference [24] Emergent Abilities of Large Language Models, does not support the sentence that a provider designated successor need not be behaviorally equivalent to its predecessor.}

\textbf{Author response:}

\begin{itemize}

\item We agree that the cited reference did not support the statement.

\item The citation has therefore been removed rather than replaced with an unrelated reference.

\end{itemize}

\textbf{Author action:}

\begin{itemize}

\item The unsupported citation was removed from the relevant sentence.

\item The now-unused bibliography entry was also deleted.

\end{itemize}

\vspace{12pt}
\concern{Reviewer \#2, Concern \#1}

\concern{Reviewer comment:}

\rvcomment{1. Please justify the construction of the primary excess-flip endpoint. F\_within is computed from all three unordered replicate pairs within each company-date, whereas F\_cross compares matched company-date-replicate cells. Please explain why the former is the appropriate noise baseline for the latter and, if possible, add a sensitivity analysis using an alternative matched or repeated-call baseline.}

\textbf{Author response:}

\begin{itemize}

\item This overlaps with Reviewer \#1, Concern \#1.

\item We addressed it by estimating same-seed repeat-call disagreement on a stratified subgrid, rather than relying only on an interpretive defense.

\item Applying that repeat-call floor as an alternative noise-floor estimate to the full-grid cross-version disagreement rate leaves the endpoint positive for both migrations: 0.189 for OpenAI and 0.375 for Gemini.

\end{itemize}

\textbf{Author action:}

\begin{itemize}

\item Section III-C was added.

\item It now:

\begin{itemize}

\item explains why the preregistered within-version floor answers the intended operational question;

\item explicitly acknowledges the different coupling structures; and

\item reports the requested same-seed repeat-call sensitivity analysis, stating that its floor is applied to the full-grid cross-version disagreement rate.

\end{itemize}

\end{itemize}

\vspace{12pt}
\concern{Reviewer \#2, Concern \#2}

\concern{Reviewer comment:}

\rvcomment{2. Please make the Gemini migration provenance independently verifiable. Table 1 gives gemini-2.5-flash -> gemini-3.5-flash with an announced shutdown of October 16, 2026, sourced to [34] as accessed on July 4, 2026. Please include an archived capture or other stable provider evidence establishing the shutdown date and designated successor as of the protocol-freeze date. If the pair was selected experimentally rather than explicitly designated by the provider, the ``vendor-declared migration'' terminology should be revised throughout.}

\textbf{Author response:}

\begin{itemize}

\item We located an archived Google deprecation-page capture dated 1 July 2026.

\item It confirms:

\begin{itemize}

\item gemini-2.5-flash;

\item shutdown date 16 October 2026; and

\item gemini-3.5-flash as the recommended replacement.

\end{itemize}

\item The ``vendor-declared'' terminology is therefore retained: the manuscript describes the pairs as vendor-declared model-version migrations.

\end{itemize}

\textbf{Author action:}

\begin{itemize}

\item Section IV-D, ``Models and Migration Pairs'' was revised.

\item We:

\begin{itemize}

\item added the archived provider evidence;

\item stated the shutdown date and designated successor explicitly;

\item explained why an archived capture is used instead of relying on a revisable live page; and

\item added the archived source as a numbered reference.

\end{itemize}

\end{itemize}

\vspace{12pt}
\concern{Reviewer \#2, Concern \#3}

\concern{Reviewer comment:}

\rvcomment{3. Please cite the SR 26-2 provisions with greater precision. Section VII states that SR 26-2 replaced SR 11-7 in April 2026 and expressly excludes generative and agentic AI from scope; this premise carries the governance-gap argument. Please identify the specific provisions supporting these statements and reconcile Table 5, which maps only SR 11-7 provisions, with the text's position that SR 26-2 is the operative framework.}

\textbf{Author response:}

\begin{itemize}

\item We agree that the governance discussion required more precise sourcing.

\item The revised manuscript now identifies:

\begin{itemize}

\item Footnote 3 for the generative/agentic AI scope carve-out;

\item Section III for aggregate model risk; and

\item Section VII for vendor and third-party models.

\end{itemize}

\item We also clarify that the carve-out delegates governance to institutional risk-management practices rather than implying an absence of governance obligations.

\end{itemize}

\textbf{Author action:}

\begin{itemize}

\item Section VII, ``Governance: SR 11-7, SR 26-2, and a Pre-Migration Control'' was substantially revised.

\item Table 5 now carries a separate column for each framework, so the mapping no longer rests on SR 11-7 alone, and each cell gives the section number together with the provision title as the guidance itself words it, rather than a section number the reader would have to look up. The five rows are the governance concepts the paper relies on: vendor and third-party responsibility; ongoing monitoring; effective challenge and validation before use; model inventory and change control; and aggregate model risk. Two further columns state, for each concept, how version migration breaks it and what control we propose.

\item The caption states that the SR 26-2 column gives the provisions that would otherwise apply, because SR 26-2's own footnote 3, attached to its Purpose and Scope section, provides that generative and agentic AI models ``are not within the scope of this guidance.'' Section VII now quotes the scope exclusion and incorporates the accompanying governance language, under which an institution's own risk-management and governance practices determine the controls for systems not covered. The carve-out therefore delegates the governance of these systems to institutional discretion rather than leaving it unaddressed, and the proposed control is framed as a way to discharge that delegated responsibility.

\item Reference [3] now cites an archived capture of SR 11-7, since the issuing page was withdrawn after supersession, and names the provisions relied upon. Reference [4] does the same for SR 26-2.

\item The governance language was narrowed to describe a governance challenge rather than an absence of regulatory obligations.

\end{itemize}

\vspace{12pt}
\concern{Reviewer \#2, Concern \#4}

\concern{Reviewer comment:}

\rvcomment{4. Please correct the description of the sector-restricted analysis. It is introduced with ``The principal pattern remains,'' but the OpenAI hit rate improves 0.539 -> 0.564 across all sectors and declines 0.590 -> 0.562 when Financials and Real Estate are excluded. Because the restricted subset is where the Altman Z`` reference is more appropriate, this reversal should be stated directly, and the Discussion's description of ''modest numerical improvement`` qualified accordingly. Please also report balanced accuracy and macro-F1 for the restricted analysis, since the manuscript argues these are required for imbalanced classes.}

\textbf{Author response:}

\begin{itemize}

\item We agree that the earlier text incorrectly described the OpenAI restricted-sample result as a continuation rather than a reversal.

\item The revised results now state explicitly:

\begin{itemize}

\item Full-sample OpenAI hit rate: 0.539 to 0.564

\item Restricted-sample OpenAI hit rate: 0.590 to 0.562

\end{itemize}

\item We also added balanced accuracy and macro-F1 for the restricted sample.

\end{itemize}

\textbf{Author action:}

\begin{itemize}

\item Section V-E, ``Secondary Accuracy Analysis'' was revised.

\item We:

\begin{itemize}

\item explicitly report the OpenAI hit-rate reversal;

\item added restricted-sample balanced accuracy and macro-F1 for both providers;

\item identify the sector-restricted sample as the more informative accuracy view; and

\item updated the Discussion and Conclusion so they no longer overstate the OpenAI improvement.

\end{itemize}

\end{itemize}

\vspace{12pt}
\concern{Reviewer \#2, Concern \#5}

\concern{Reviewer comment:}

\rvcomment{5. Please report uncertainty for the aggregate secondary comparisons or qualify them. Balanced accuracy changes 0.538 -> 0.570 for OpenAI and 0.488 -> 0.408 for Gemini, and macro-F1 changes 0.517 -> 0.565 and 0.463 -> 0.325, without intervals, unlike the flip-conditional analysis. Confidence intervals or paired bootstrap estimates would clarify whether these are inferential or descriptive.}

\textbf{Author response:}

\begin{itemize}

\item We added ticker-clustered bootstrap confidence intervals for the changes in balanced accuracy and macro-F1.

\item OpenAI:

\begin{itemize}

\item Balanced accuracy: +0.032 [-0.005, 0.070]

\item Macro-F1: +0.048 [0.007, 0.090]

\end{itemize}

\item Gemini:

\begin{itemize}

\item Balanced accuracy: -0.080 [-0.140, -0.028]

\item Macro-F1: -0.138 [-0.203, -0.078]

\end{itemize}

\item The OpenAI improvement is therefore characterized as weak, while the Gemini deterioration is clearly resolved.

\end{itemize}

\textbf{Author action:}

\begin{itemize}

\item Section V-E now reports bootstrap confidence intervals for both secondary metrics and both migration pairs.

\item The corresponding claims in the Results, Discussion, and Conclusion were moderated to match the inferential evidence.

\end{itemize}

\vspace{12pt}
\concern{Reviewer \#2, Concern \#6}

\concern{Reviewer comment:}

\rvcomment{6. Please reconcile the temperature subgrid with the primary result. The T = 0 subgrid estimates are 0.153 for OpenAI and 0.222 for Gemini, against full-grid values of 0.112 and 0.381 at the same nominal setting. The wide intervals are explained, but the difference in point estimates is not.}

\textbf{Author response:}

\begin{itemize}

\item We verified that the difference between the full-grid and subgrid T=0 estimates is due to cell selection, not temperature.

\item Re-running the primary estimator on exactly the 24 subgrid cells reproduces:

\begin{itemize}

\item OpenAI: 0.153

\item Gemini: 0.222

\end{itemize}

\end{itemize}

\textbf{Author action:}

\begin{itemize}

\item Section V-F, ``Temperature Sensitivity'' was revised.

\item We:

\begin{itemize}

\item added the restricted-cell recomputation;

\item explain that the T=0 row applies the primary estimator to the small stratified subset; and

\item clarify that it serves as the anchor for comparison with T=0.7 and T=1.0, rather than as an independent full-grid estimate.

\end{itemize}

\end{itemize}

\vspace{12pt}
\concern{Reviewer \#2, Concern \#7}

\concern{Reviewer comment:}

\rvcomment{7. Please moderate the ``decision boundary'' interpretation. That cells unstable within the incumbent version show higher cross-version flip rates is a useful observation, but within-version instability does not establish proximity to an underlying decision boundary. Describing these as cases with greater pre-existing output instability would be sufficient.}

\textbf{Author response:}

\begin{itemize}

\item We agree that within-version instability does not establish proximity to an underlying decision boundary.

\item We removed that interpretation.

\end{itemize}

\textbf{Author action:}

\begin{itemize}

\item Section V-D, ``Robustness of Cross-Version Agreement'' was revised.

\item The text now states only that migration changes are concentrated among observations with greater pre-existing within-version output instability.

\item It explicitly states that proximity to an underlying decision boundary is not measured.

\end{itemize}

\vspace{12pt}
\concern{Reviewer \#2, Concern \#8}

\concern{Reviewer comment:}

\rvcomment{8. Please clarify the noise-floor interpretation. A positive cluster-bootstrap interval shows that the defined cross-version disagreement exceeds the defined within-version baseline. The statement that the result cannot be explained ``solely'' by residual sampling variability should be justified under the estimator's assumptions or phrased more narrowly.}

\textbf{Author response:}

\begin{itemize}

\item We agree that the previous language was too broad.

\item The result establishes that the defined cross-version disagreement exceeds the defined within-version baseline under the estimator's assumptions.

\item It does not exclude every possible sampling-related explanation.

\end{itemize}

\textbf{Author action:}

\begin{itemize}

\item Section V-A, ``Migration-Induced Decision Instability'' was revised.

\item The previous ``cannot be explained solely by sampling variability'' language was replaced with a narrower statement describing exactly what the confidence interval establishes.

\end{itemize}

\vspace{12pt}
\concern{Reviewer \#2, Concern \#9}

\concern{Reviewer comment:}

\rvcomment{9. Please retain the current limitations on systemic-risk interpretation. The manuscript appropriately states that it demonstrates a potential shared-update mechanism rather than realized cross-institutional correlation or losses, and this distinction should remain explicit in the Abstract, Discussion, and Conclusion.}

\textbf{Author response:}

\begin{itemize}

\item We agree and retained this limitation throughout the manuscript.

\item The study demonstrates a deployment-level shared-update mechanism, not realized cross-institutional correlation, losses, or market-level systemic effects.

\end{itemize}

\textbf{Author action:}

\begin{itemize}

\item The limitation was verified and retained in the Abstract and Discussion.

\item Section IX, Conclusion was strengthened to state explicitly that the study does not estimate:

\begin{itemize}

\item typical migration effects;

\item realized cross-institutional correlation; or

\item market-level losses.

\end{itemize}

\end{itemize}

\end{document}
